\documentclass[11pt]{article}
\usepackage[a4paper,margin=1.45cm]{geometry}
\usepackage{fontspec}
\usepackage{xeCJK}
\setmainfont{TeXGyreTermes}
\setCJKmainfont{Noto Serif CJK JP}
\setsansfont{TeXGyreHeros}
\setCJKsansfont{Noto Sans CJK JP}
\usepackage{amsmath,amssymb,mathtools,array,booktabs}
\usepackage[protrusion=true,expansion=false]{microtype}
\setlength{\parindent}{1.0em}
\setlength{\parskip}{0pt}
\allowdisplaybreaks
\setlength{\emergencystretch}{30em}
\sloppy
\hbadness=10000
\vbadness=10000
\hfuzz=10pt
\vfuzz=10pt
\raggedbottom
\providecommand{\mG}{\mathfrak{G}}
\providecommand{\mA}{\mathfrak{A}}
\providecommand{\mB}{\mathfrak{B}}
\providecommand{\mC}{\mathfrak{C}}
\providecommand{\mD}{\mathfrak{D}}
\providecommand{\mE}{\mathfrak{E}}
\providecommand{\mH}{\mathfrak{H}}
\providecommand{\mK}{\mathfrak{K}}
\providecommand{\mM}{\mathfrak{M}}
\providecommand{\mN}{\mathfrak{N}}
\providecommand{\mR}{\mathfrak{R}}
\providecommand{\mS}{\mathfrak{S}}
\providecommand{\mT}{\mathfrak{T}}
\providecommand{\mU}{\mathfrak{U}}
\providecommand{\mO}{\mathfrak{O}}
\providecommand{\mo}{\mathfrak{o}}
\providecommand{\ma}{\mathfrak{a}}
\providecommand{\mb}{\mathfrak{b}}
\providecommand{\mc}{\mathfrak{c}}
\providecommand{\md}{\mathfrak{d}}
\providecommand{\mf}{\mathfrak{f}}
\providecommand{\mh}{\mathfrak{h}}
\providecommand{\ml}{\mathfrak{l}}
\providecommand{\mn}{\mathfrak{n}}
\providecommand{\mr}{\mathfrak{r}}
\providecommand{\mt}{\mathfrak{t}}
\providecommand{\mx}{\mathfrak{x}}
\providecommand{\mz}{\mathfrak{z}}
\providecommand{\mZ}{\mathfrak{Z}}
\providecommand{\mL}{\mathfrak{L}}
\providecommand{\iso}{\simeq}
\providecommand{\ann}{\operatorname{Ann}}
\providecommand{\tuple}[1]{(#1)}
\providecommand{\mat}[1]{\begin{pmatrix}#1\end{pmatrix}}
\begin{document}

\setcounter{footnote}{0}
\section*{34. 超複素量と表現論}
\emph{続き：第I章 §§5--7 および第II章 §§8--14。}

\subsection*{§ 5. 完全可約群}

群は、有限個の単純群の直積であるとき、\emph{完全可約}であるという：
\[
        \mG=\mA_1\times\cdots\times\mA_n.
\]
この場合、群
\[
        \mA_1\times\cdots\times\mA_n
        \supset \mA_1\times\cdots\times\mA_{n-1}
        \supset\cdots\supset \mA_1\supset \mE
\]
は組成列をなす。その長さは $n$ であり、組成因子は
\[
        \sim \mA_n,\ \mA_{n-1},\ldots,\mA_1
\]
である。

\emph{$\mG$ が完全可約なら、任意の正規因子は直因子であり、他方の因子は、$\mG$ の与えられた積分解に現れる単純群の積として選ぶことができる。正規因子 $\mH$ 自身も完全可約である。}

\emph{証明。} つぎを得る：
\begin{equation}
        \mG=\mH\mA_1\mA_2\cdots\mA_n.             \tag{3}
\end{equation}
$\mH_i=\mH\mA_1\cdots\mA_i$ とおく。このとき $\mH_{i+1}=\mH_i\mA_{i+1}$ である。あるいは $\mA_{i+1}\leq \mH_i$ であり、そのとき $\mH_{i+1}=\mH_i$ だから、因子 $\mA_{i+1}$ を (3) から省く。あるいは $\mH_i\cap\mA_{i+1}$ が $\mA_{i+1}$ の真の正規因子であり、したがって $\mE$ であって、このとき $\mH_i\times\mA_{i+1}$ は直積である。余分な因子を省いたのち、(3) は
\[
        \mG=\mH\times\mA_{i_1}\times\cdots\times\mA_{i_r}
\]
となるので、$\mH$ は直因子である。さらに
\[
        \mH\sim \mG/(\mA_{i_1}\times\cdots\times\mA_{i_r})
        \sim \mA_{j_1}\times\cdots\times\mA_{j_\mu},
\]
ただし $\mA_{j_1},\ldots,\mA_{j_\mu}$ は残りの $\mA_i$ である。ゆえに $\mH$ は完全可約である。

\emph{系。} \emph{完全可約群を直接分解不能因子に分解する任意の分解は、単純因子への分解であり、したがって組成列を与える。これから因子が同型を除いて一意に定まることが従う。}

\subsection*{§ 6. 体に関する加群。超複素系}

作用素付き完全可約群のほとんど自明な例は、必ずしも可換とは限らない体、すなわち斜体に関する有限基底をもつ加群によって与えられる。

$\mG$ を右 $K$-加群とし、$K$ の単位元 $e$ は同時に恒等作用素であるとする：$\mG$ の $a$ について $ae=a$。ある $a$ から作られる任意の部分加群 $aK$ は単純であり、すなわちその任意の非零元から作られる加群に等しい。したがって、$\mH$ がいくつかの元から作られる部分加群で、$a$ が $\mH$ に属さない元であるなら、$\mH\cap aK=\mE$ であり、$\mH+aK$ は直和である。そこで任意の基底元 $a_1$ から始め、新しい基底元 $a_2,a_3,\ldots$ を順に付け加えることができ、$\mG$ は直和
\[
        \mG=a_1K+a_2K+\cdots+a_nK
\]
として得られる。\emph{したがって $\mG$ は完全可約である。}

数 $n$、すなわち組成列の長さを、$K$ に関する $\mG$ の\emph{階数}という。$\mG$ の元の表示の一意性により（また $e$ が恒等作用素であると仮定したことにより）、基底 $a_1,\ldots,a_n$ は線型独立である。

任意の部分加群 $\mH$ は直和因子である：
\[
        \mG=\mH+\mR=(h_1K+\cdots+h_sK)+(r_1K+\cdots+r_{n-s}K),
\]
すなわち、\emph{$\mH$ の任意の線型独立な基底は、$\mG$ の線型独立な基底へ拡張できる。} しかも $r_i$ は、もとの基底元 $a_i$ の中から選ぶこともできる。

$(a_1,\ldots,a_n)$ と $(c_1,\ldots,c_n)$ を線型独立な基底とする。このとき
\[
        c_k=\sum_i a_i\pi_{ik},
\]
すなわち行列で書けば、
\[
        (c_1,\ldots,c_n)=(a_1,\ldots,a_n)P,
        \qquad
        P=\begin{pmatrix}
        \pi_{11}&\cdots&\pi_{1n}\\
        \vdots&&\vdots\\
        \pi_{n1}&\cdots&\pi_{nn}
        \end{pmatrix}.
\]
逆に、
\[
        (a_1,\ldots,a_n)=(c_1,\ldots,c_n)Q.
\]
したがって
\[
        (a_1,\ldots,a_n)=(a_1,\ldots,a_n)PQ,
        \qquad
        (c_1,\ldots,c_n)=(c_1,\ldots,c_n)QP,
\]
であり、$a_i$ と $c_i$ が線型独立であるから、
\[
        PQ=QP=E.
\]

同時に環でもある特殊な $K$-加群が「超複素系」である。

環 $\mo$ が同時に可換体 $K$ に関する右加群であるとき、それが次を満たせば、$K$ に関する\emph{超複素系}（文献では「$K$ 上の代数」とも呼ばれる）という：
\begin{enumerate}
\item 階数が有限である（線型独立な基底を $u_1,\ldots,u_n$ とする）。
\item $ab\cdot x=a\cdot bx=ax\cdot b$。これは、$K$ が $\mo$ と可換に結合されている、と言い表される。
\item $K$ の単位元 $e$ は同時に恒等作用素である：$\mo$ の $a$ について $ae=a$。
\end{enumerate}

\[
        \Bigl(\sum_i u_i\alpha_i\Bigr)\Bigl(\sum_k u_k\beta_k\Bigr)
        =\sum_{i,k}(u_i u_k)(\alpha_i\beta_k)
\]
であるから、この系は、$K$ に加えて\emph{乗法表}が与えられると一意に定まる。すなわち各積 $u_i u_k$ が $u_j$ によってどのように表されるか、
\[
        u_i u_k=\sum_j u_j\gamma^j_{ik}
\]
が知られればよい。$\gamma^j_{ik}$ は、結合法則から従うよく知られた関係式を満たさなければならない。

$\mo$ が、しばしば仮定するように、単位元 $e$ をもつなら（すべての $a$ について $ea=ae=a$）、$K$ の元 $x$ は $ex$ と同一視でき、$K$ は $\mo$ の部分体とみなせる。そのとき超複素系は、\emph{中心に含まれる体に関して有限階数をもつ環}とも記述できる。

超複素系の例は、有限群の\emph{群環}である。群の元を基底元 $u_i$ とし、群の乗法を乗法として用いる。体 $K$ は任意である。

\footnotetext{条件 2 は、$K$ の元による乗法が、$\mo$ を $\mo$-左加群としても $\mo$-右加群としても見たときの $\mo$ の準同型である、ということである。3 によって、これらの準同型はさらに同型である。}

\subsection*{§ 7. 行列}

環 $\mo$ の元 $\pi_{ik}$ を成分にもつ正方行列 $(\pi_{ik})$ は、通常の行列積 $\sum_j\pi_{ij}\rho_{jk}=\sigma_{ik}$ と加法 $\pi_{ik}+\rho_{ik}=\sigma_{ik}$ によって環をなす。

体 $K$ の元を成分とする行列で、右逆元と左逆元の両方が存在するものを\emph{正則}という。

§6 で見たように、\emph{$K$-加群の二つの線型独立な基底の間の変換行列は正則である。}

\emph{正則性には右逆元で十分である。} 実際、線型独立な基底 $(a_1,\ldots,a_n)$ をもつ右 $K$-加群（不定元 $a_1,\ldots,a_n$ における線型形式の加群）を作り、
\[
        (c_1,\ldots,c_n)=(a_1,\ldots,a_n)P
\]
とおく。このとき
\[
        (c_1,\ldots,c_n)P^{-1}=(a_1,\ldots,a_n)PP^{-1}=(a_1,\ldots,a_n),
\]
であるから、$(c_1,\ldots,c_n)$ は再び基底であり、したがって再び線型独立であり、行列 $P$ は正則である。さらに右逆元は左逆元に等しい。

同様に、左逆元で十分である。

\emph{$P$ が左逆元をもつなら、$P$ は左零因子ではない。}
\emph{証明。} $PA=0$ から $A=P^{-1}PA=0$ が従う。

したがって $P$ は右逆元もただ一つしかもたず、それは §6 によって左逆元と一致する。

通常どおり $P_x$ が行列 $(\pi_{ik}x)$ を表し、$C_{ik}$ が $ik$ の位置に単位元をもち他のすべての位置に零をもつ行列を表すなら、
\[
        (\pi_{ik})=\sum_{i,k} C_{ik}\pi_{ik};
\]
したがって $K$ 上の行列は階数 $n^2$ の $K$-加群をなす。$C_{ik}$ は関係式
\[
        C_{ij}C_{jk}=C_{ik},
        \qquad
        C_{ij}C_{\bar j k}=0,\quad j\ne\bar j
\]
を満たす。したがって $K$ 上の行列は有限 $K$-加群をなし、$K$ が可換なら超複素系をなす。

\section*{第II章　非可換イデアル論}

\subsection*{§ 8. 環に対する準同型定理}

$\ma$ が $\mo$ の両側イデアルなら、剰余類領域 $\mo/\ma$ は、容易にわかるように $\mo$-加群であるだけでなく環でもある。$\mo$ の任意の準同型像は再び環であり、両側イデアル $\ma$ による剰余類環 $\mo/\ma$ と同型である。証明は群の場合と同じである。

とくに、$\mo$ が体であれば、$(0)$ と $\mo$ 以外のイデアルは存在しない。したがってこの場合の準同型は、同型であるか、すべての元を零へ写すかのいずれかである。

とくに、環 $\mo$ が右 $K$-加群であり、$K$ は $\mo$ と元ごとに可換な環であると仮定する：
\[
        ab\cdot x=a\cdot bx=ax\cdot b
\]
（たとえば、可換体 $K$ に関する超複素系の場合）。このとき、許容される右、左、または両側イデアルとしては、同時に $K$-加群であるものだけを考え、環準同型に加えて、$K$ による乗法に関する作用素準同型も要求する。許容イデアルは $\mo$ と $K$ に関する双加群になる。（超複素系では、単に「イデアル」といえば、基礎にある係数体 $K$ に関して許容されるものだけを常に意味する。）この拡張された意味でも、上の準同型定理は有効である。

環準同型の一例は、乗数領域 $\mo$ から絶対乗数領域への移行である（§1、例4）。したがって絶対乗数領域は剰余類環 $\mo/\md$ と同型であり、ここで $\md$ は $\mM$ を消す $\mo$ の元からなる両側イデアルである。

準同型定理が成り立つことから、§2 と同様に、環同型および作用素同型について第一および第二同型定理が従う。

$\mo$ の部分集合のすべての積 $\mA\mB$ は、加群積として理解する。すなわち $\mA\mB$ は、$a$ が $\mA$ に、$b$ が $\mB$ に属するすべての和 $\sum ab$ の全体であり、$a\mB$ は $b$ が $\mB$ に属するすべての和 $\sum ab$ の全体である。$\mB$ が、ある環を作用素領域とする右加群なら、積 $\mA\mB$、それぞれ $a\mB$ もまた右加群である。とくに、右イデアル、係数領域 $K$ に関して許容されるイデアルなどである。$(\mA,\mB)$ が加群 $\mA$ と $\mB$ の和を表すなら、計算法則は
\[
        \mA\mB\cdot\mC=\mA\cdot\mB\mC,
\]
\[
        \mA(\mB,\mC)=(\mA\mB,\mA\mC),
        \qquad
        (\mA,\mB)\mC=(\mA\mC,\mB\mC)
\]
である。直和は $\mA+\mB$ と書く（§4）。

以下では、右イデアルについて極大条件または極小条件（§3）を満たす環をしばしば考える。とくに超複素系はこれに属する。いま述べた約束により、イデアルとしては $K$-加群だけが許容され、その階数が有界に保たれるからである（§6）。

\subsection*{§ 9. 冪等元。右イデアルへの直和分解}

$\mo$ を単位元をもつ環とする。

$\mr$ が右イデアルなら、$\mr\mo\subseteq\mr$ であり、単位元の存在によって実は $\mr\mo=\mr$ である。左イデアルについては対応して $\mo\ml=\ml$ である。

元 $c$ は $c^2=c$ のとき冪等元と呼ばれる。

\emph{$\mo=\mr_1+\cdots+\mr_n$ が右イデアルへの分解であり、かつ}
\[
        e=e_1+\cdots+e_n,
        \qquad e_i\in\mr_i,
\]
\emph{なら、}
\[
        e_i^2=e_i,
        \qquad e_i e_k=0\quad(i\ne k)
        \quad\hbox{（「直交関係」）},
\]
\[
        \mr_i=e_i\mo
\]
\emph{である。}

\emph{証明。} $r$ を $\mr_1$ の元とする。このとき
\[
        r=er=e_1r+\cdots+e_nr,
        \qquad e_ir\in\mr_i,
\]
であるが、一方で
\[
        r=r+0+\cdots+0
\]
でもある。したがって
\[
        e_1r=r,
        \qquad e_ir=0\quad(i\ne1).
\]
第一式から $\mr_1\subseteq e_1\mo$ が従う。他方 $e_1\mo\subseteq\mr_1$ であるから、$e_1\mo=\mr_1$。$r=e_1$ と特殊化すると、
\[
        e_1^2=e_1,
        \qquad e_i e_1=0\quad(i\ne1)
\]
を得る。添字 $1$ を $k$ で置き換えても同じである。

\emph{逆に、$e=\sum e_i$、$e_i e_k=0$ $(i\ne k)$、$e_i^2=e_i$ であり、$\mr_i=e_i\mo$ とおくなら、$\mo=\mr_1+\cdots+\mr_n$ である。}

\emph{証明。} $\mo$ の任意の元 $r$ は
\[
        r=er=e_1r+\cdots+e_nr
\]
であるから、
\[
        \mo=(e_1\mo,\ldots,e_n\mo).
\]
しかし表示は一意である。なぜなら
\[
        0=e_1a_1+\cdots+e_na_n
\]
から、$e_i$ を掛けると
\[
        e_i^2a_i=e_i a_i=0
\]
が従うからである。

右側表示
\[
        \mo=\mr_1+\cdots+\mr_n=e_1\mo+\cdots+e_n\mo
\]
から、それに応じて左側表示
\[
        \mo=\ml_1+\cdots+\ml_n=\mo e_1+\cdots+\mo e_n
\]
が従う。

$\mr_i$ が直接分解不能なら、$\ml_i$ もそうである。たとえば $\ml_1$ の分解があれば、
\[
        \ml_1=\mo e_1=\mo e_1'+\mo e_1'',
        \qquad
        \mo=\mo e_1'+\mo e_1''+\mo e_2+\cdots+\mo e_n
\]
を意味する。このことから右分解
\[
        \mo=e_1'\mo+e_1''\mo+e_2\mo+\cdots+e_n\mo
        =\mr_1'+\mr_1''+\mr_2+\cdots+\mr_n,
\]
\[
        \mr_1\cong \mo/(\mr_2+\cdots+\mr_n)=\mr_1'+\mr_1''
\]
が生じ、$\mr_1$ は分解可能になってしまう。

\footnotetext{もう一つの逆もある。$e_i e_k=0$、$e_i^2=e_i$、$c=\sum e_i$ なら、$c$ は環 $e_1\mo+\cdots+e_n\mo$ の左単位元である。和が直和であることは、上と全く同じ方法でわかる。}

\subsection*{§ 10. 両側イデアルへの分解}

再び $\mo$ を単位元をもつ環とし、
\begin{equation}
        \mo=\ma_1+\cdots+\ma_n                                  \tag{1}
\end{equation}
を $\mo$ の両側イデアルへの分解とする。このとき直交関係はイデアルについても成り立つ：
\begin{equation}
        \ma_i\ma_k=0\quad(i\ne k),
        \qquad
        \ma_i^2=\ma_i.                                             \tag{2}
\end{equation}

\emph{証明。} $\mb_1=\ma_2+\cdots+\ma_n$、$\mo=\ma_1+\mb_1$ とする。このとき $\ma_1\mb_1\leq\ma_1\cap\mb_1=0$ であるから $\ma_1\mb_1=0$、したがってなおさら
\[
        \ma_1\ma_i=0\quad(i\ne1);
\]
また
\[
        \ma_1=\ma_1\mo=\ma_1(\ma_1+\mb_1)=(\ma_1^2,\ma_1\mb_1)=\ma_1^2.
\]

逆に、\emph{(1) と (2) から、加群 $\ma_i$ は両側イデアルであることが従う。}
\emph{証明。}
\[
        \mo\ma_i=(\ma_1+\cdots+\ma_n)\ma_i=\ma_i^2=\ma_i,
        \qquad
        \ma_i\mo=\ma_i.
\]

\emph{環 $\ma_i$ における右イデアルは、$\mo$ における右イデアルである。}
\emph{証明。}
\[
        \mr\mo=\mr(\ma_i+\mb_i)=(\mr\ma_i,\mr\mb_i)=(\mr,0)=\mr.
\]

\emph{$\mr$ が $\mo$ の右イデアルなら、$\mr$ は $\ma_i$ に含まれる右イデアル $\mr\ma_i$ の直和である。}
\emph{証明。} $\mr=\mr\mo=(\mr\ma_1,\ldots,\mr\ma_n)$ であり、$\mr\ma_i\mo=\mr\ma_i$ だから、$\mr\ma_i$ は $\mr$ に含まれる右イデアルである。$\mr\ma_i\leq\ma_i$ なので、和は直和である：
\[
        \mr=\mr\ma_1+\cdots+\mr\ma_n.
\]

\emph{とくに、$\mr$ が直接分解不能なら、$\mr$ は一つの成分しかもつことができない。したがって $\mr$ はある $\ma_i$ に含まれなければならない。}

これらの定理によって、個々の $\ma_i$ のイデアル論がわかれば、$\mo$ におけるイデアル論は制御される。したがって以下では、ほとんどの場合、両側分解不能な環に制限する。

\emph{異なる $\ma_i$ に含まれる二つの右イデアルは、決して作用素同型ではありえない。}
\emph{証明。} $\ma_1$ に含まれるイデアルは、他のすべての $\ma_k$ によって消されるが、$\ma_1$ によっては消されない。これに対して、$\ma_2$ に含まれるイデアルは $\ma_1$ によって消される。

したがって、環 $\mo$ を両側分解不能イデアル $\ma_1+\cdots+\ma_n$ に分解し、各 $\ma_i$ をさらに片側の直接分解不能な右イデアル $\mr_i'+\mr_i''+\cdots$ に分解できるなら、作用素同型な $\mr$ は同じ $\ma_i$ に属するものの中で探さなければならない。しかし、同じ $\ma_i$ の中にもなお異なる $\mr$ の類、すなわち作用素同型でない類が存在することがある。次の例がそれを示す。

$K$ を有理数体とし、
\[
        \mo=e_1K+e_2K+uK
\]
を超複素系とする。$e_i$ と $u$ は $K$ の数と可換であり、乗法表は
\[
\begin{array}{c|ccc}
     & e_1&e_2&u\\ \hline
 e_1 & e_1&0&u\\
 e_2 & 0&e_2&0\\
 u   & 0&u&0
\end{array}
\]
である。単位元は $e=e_1+e_2$ である。$e_i$ は直交関係を満たす。したがって右分解は
\[
        \mo=e_1\mo+e_2\mo=(e_1,u)+(e_2),
\]
左分解は
\[
        \mo=\mo e_1+\mo e_2=(e_1)+(e_2,u)
\]
である。$e_2\mo$ と $\mo e_1$ は見て明らかに分解不能なので、$\mo e_2$ と $e_1\mo$ も分解不能でなければならない。

両側分解は不可能である。というのは、その場合、一方の成分は $e_1\mo$ を含み、他方は $e_2\mo$ を含まなければならないが、前者は $u$ を含み、後者も $ue_2=u$ を含むことになるからである。

しかし $e_1\mo$ と $e_2\mo$ は作用素同型ではない。$K$ に関する階数が異なるからである。

\emph{$\mo=\ma_1+\cdots+\ma_n$ とし、$\ma_i$ を両側分解不能とする。このとき $\ma_i$ は一意に定まる。}

\emph{証明。} さらに $\mo=\mc_1+\cdots+\mc_n$ と仮定する。このとき
\[
        \ma_i=\ma_i\mo=(\ma_i\mc_1,\ldots,\ma_i\mc_n).
\]
この後の和は、$\ma_i\mc_k\leq\mc_k$ かつ $\leq\ma_i$ であるため直和である。しかし $\ma_i$ は分解不能なので、$\ma_i\mc_k$ は一つ、$\ma_i\mc_{j_i}$ を除いてすべて零でなければならない。したがって
\[
        \ma_i=\ma_i\mc_{j_i}.
\]
同様に、
\[
        \mc_{j_i}=\mo\mc_{j_i}=\ma_1\mc_{j_i}+\cdots+\ma_n\mc_{j_i},
\]
であり、$\ma_i\mc_{j_i}\ne0$ であるから、残りの項はすべて $0$ である。ゆえに
\[
        \mc_{j_i}=\ma_i\mc_{j_i}=\ma_i,
\]
q.e.d.

\subsection*{§ 11. 中心}

環 $\mo$ の中心 $\mZ$ とは、すべての環元と可換なすべての $z$ の全体である（すべての $a$ について $za=az$）。$\mZ$ は可換環である。

$\mo$ の任意の片側イデアル $\ma$ は、$\mZ$ の中に「収縮イデアル」$\ma\cap\mZ$ をもつ。$\ma$ と $\mZ$ はどちらも $\mZ$-加群であるから、$\ma\cap\mZ$ もそうである。

$\mZ$ の任意のイデアル $\mA$ は拡張イデアルをもつ。すなわち $\mA$ によって $\mo$ の中に生成されるイデアル $\ma=(\mA,\mo\mA)$ である。これは両側である。なぜなら
\[
        \ma\mo=(\mo\mA,\mA)\mo=(\mo\mA\mo,\mA\mo)
        = (\mo^2\mA,\mo\mA)\subseteq \mo\mA\subseteq\ma
\]
だからである。$\mo$ が単位元をもつと仮定すれば、$\ma=\mo\mA$ と書ける。この場合、次の定理が成り立つ。

\emph{$\mo$ の任意の両側分解は、中心の分解と一対一に対応する：}
\begin{equation}
        \mo=\ma_1+\cdots+\ma_n,
        \qquad \mA_i=\ma_i\cap\mZ                         \tag{1}
\end{equation}
\emph{から}
\begin{equation}
        \mZ=\mA_1+\cdots+\mA_n,
        \qquad \mo\mA_i=\ma_i,                              \tag{2}
\end{equation}
\emph{が従い、また逆も成り立つ。}

\emph{証明。} $z$ を中心元とし、
\begin{equation}
        z=z_1+\cdots+z_n                                    \tag{3}
\end{equation}
を (1) に従うその分解とする。このとき $z_i$ は中心元である。実際、
\[
        za=z_1a+\cdots+z_na=az=az_1+\cdots+az_n.
\]
和が直和であり、かつ $\ma_i$ が両側であることから、
\[
        z_i a=az_i
\]
が従う。したがって $z_i$ は $\ma_i\cap\mZ=\mA_i$ に属し、(3) は主張された和分解を与える。とくに $e=\sum e_i$ であるから、$e_i$ は中心元である。最後に
\[
        z=ze_1+\cdots+ze_n,
\]
したがって
\[
        \mA_i=\mZ e_i,
        \qquad
        \mo\mA_i=\mo\mZ e_i=\mo e_i=\ma_i.
\]

逆に、$\mZ=\mA_1+\cdots+\mA_n$ が中心の分解であるなら、§9 により
\[
        \mA_i=\mZ e_i,
        \qquad e_i^2=e_i,
        \qquad e_i e_k=0\quad(i\ne k)
\]
である。したがって、直交関係 \hbox{§9} により、
\[
\begin{aligned}
        \mo&=\mo e=\mo e_1+\cdots+\mo e_n
        =\mo\mZ e_1+\cdots+\mo\mZ e_n\\
        &=\mo\mA_1+\cdots+\mo\mA_n
        =\ma_1+\cdots+\ma_n .
\end{aligned}
\]
いま $\ma_i\cap\mZ=\mA_i'$ とすると、主張の前半から
\[
        \mZ=\mA_1'+\cdots+\mA_n'\quad\hbox{が直和である}
\]
ことがわかる。しかし $\mA_i'\supseteq\mA_i$ であるから、§4, 2 により $\mA_i'=\mA_i$ である。

\emph{系。} $\ma_i$ と $\mA_i$ は、同時に両側直接分解可能または分解不能である。

\subsection*{§ 12. 冪零イデアル}

イデアル $\mc$ は $\mc^\rho=0$ となるとき\emph{冪零}であるという。

\emph{例。} §10 のイデアル $(u)$。

\emph{$\mo$ に冪零右イデアル $\mr$ が存在するなら、冪零左イデアル、実際には両側イデアルも存在する。}

\emph{証明。} $\mr^\rho=0$ とする。このとき $\mo\mr$、あるいは $\mo$ が単位元を含まず $\mo\mr$ が零になりうるなら $(\mo\mr,\mr)$ は左イデアルであり、
\[
        (\mo\mr)^\rho=\mo\mr\mo\mr\cdots\mo\mr\subseteq\mo\mr\mr\cdots\mr=\mo\mr^\rho=0.
\]

\emph{二つの冪零右イデアルの和は、再び冪零右イデアルである。}

\emph{証明。} $\mc^\rho=0$、$\md^\sigma=0$ とする。
\[
        (\mc,\md)^{\rho+\sigma-1}=(\ldots,\md\mc\cdots\md\cdots\mc,\ldots)
\]
の各項には、少なくとも $\rho$ 個の因子 $\mc$、または少なくとも $\sigma$ 個の因子 $\md$ のどちらかが含まれる。前者の場合、
\[
        \md\mc\cdots\md\cdots\mc\cdots\subseteq\md\mc\mc\cdots\mc=\md\mc^\rho=0;
\]
後者の場合には $\md$ の因子について同じ議論を行う。したがって
\[
        (\mc,\md)^{\rho+\sigma-1}=0.
\]

いま $\mo$ で右イデアルについて極大条件が満たされるなら、極大な冪零右イデアルが存在する。それは他のすべての冪零右イデアルを含む。そうでなければ、そのようなイデアルとの和を作ることができるからである。それを $\mc$ と呼ぶ。$\mo\mc$ もまた冪零右イデアルであるから、$\mo\mc\subseteq\mc$ であり、したがって $\mc$ は左イデアルである。任意の冪零左イデアル $\md$ も同様に $\mc$ に含まれる。なぜなら $(\md,\md\mo)$ は冪零右イデアルだからである。したがって $\md\subseteq\mc$。\emph{こうして、他のすべての右および左の冪零イデアルを含む極大両側冪零イデアル $\mc$、すなわち「根基」が存在する。}

§10 の例では、$(u)$ が極大冪零イデアルである。

根基が零イデアルであるとき、「根基なしの環」（「半単純環」、「Dedekind 系」）という。根基を法とする剰余類環は常に根基なしの環である。

\emph{$\mZ$ が $\mo$ の中心であり、$\mc$ が $\mo$ の根基であるなら、$\mC=\mc\cap\mZ$ は $\mZ$ の根基である。}

\emph{証明。} $\mC$ が冪零であることは明らかである。もし $\mZ$ の中により大きい冪零イデアル $\overline{\mC}$ が存在するなら、それを $\mc$ に完全には含まれない冪零イデアル $\overline{\mc}$ に拡張できる。なぜなら $((\overline{\mC}\mo)^\rho=\overline{\mC}^{\rho}\mo^\rho=0)$ だからである。

\emph{帰結。} $\mo$ が根基なしの環なら、$\mZ$ もそうである。しかし逆は正しくない。§10 の例が示している。$\mo/\mc$ の中心は $\mZ/\mC$ と同型な環を含むが、これは真の部分環でありうる（§10 の例）。

\subsection*{§ 13. 完全可約環}

§5 に従って、環は有限個の単純右イデアルの直和であるとき、右完全可約であるという。

\emph{単位元をもつ右完全可約環は、$(0)$ と異なる冪零イデアルをもたず、したがって根基をもたない。}

\emph{証明。} 任意の右イデアル $\mr$ は直和因子である（§5）：
\[
        \mo=\mt+\mr=e_1\mo+e_2\mo.
\]
$e_2^2=e_2$ なので、$e_2^\rho=e_2$ でもある。もし $\mr$ が冪零なら、$e_2^\rho=0$ であり、したがって $e_2=0$、ゆえに $\mr=0$ である。q.e.d.

次に二つの逆を示す。第一に、\emph{右イデアルについて極小条件を満たす根基なしの環は、右イデアルに関して完全可約である。}

\emph{証明。} $\mt$ を $0$ でない極小右イデアルとする。$\mt$ が冪等元を含むことを示す。

$\mt^2\subseteq\mt$ かつ $\mt^2\ne0$ であるから、$\mt^2$ は右イデアルなので $\mt^2=\mt$ である。したがって、$a\mt\ne0$ となる $a$ が $\mt$ の中に存在する。このとき $a\mt$ は右イデアルだから、必ず $a\mt=\mt$ である。

$a$ によって消される $\mt$ のすべての $b$（$ab=0$）の全体は右イデアルであり、$\mt$ とは等しくないので $0$ である。したがって $ab=0$ なら $b=0$ である。

$\mt=a\mt$ であるから、$a$ は $a=ac$、$c\ne0$ と表されなければならない。すると
\[
        ac=ac^2,
        \qquad a(c-c^2)=0,
        \qquad c-c^2=0,
        \qquad c^2=c
\]
が従う。つぎに $\mt$ が直和因子であることを示す。集合 $c\mo$ は $\mt$ に含まれる非零右イデアルである。実際、$c^2=c$ は $c\mo$ に属する。したがってこれは $\mt$ である。$\mo$ の任意の元 $r$ は
\[
        r=cr+(r-cr) \qquad \hbox{（片側 Peirce 分解）}
\]
と表される。元 $cr$ はイデアル $\mt$ をなす。元 $r-cr$ は別の右イデアル $\mr$ をなし、$c$ によって消される：$c\mr=0$。したがって
\[
        \mo=(\mt,\mr).
\]
$\mt$ の元は $c$ によって消されないので、この和は直和である：
\begin{equation}
        \mo=\mt+\mr.                                             \tag{1}
\end{equation}

極小条件により $\mo$ を直接分解不能イデアルの直和
\[
        \mo=\mr_1+\cdots+\mr_n
\]
として表すと、$\mr_i$ は単純、すなわち極小でなければならない。というのは、たとえば $\mr_1$ が単純でなく、$\mt$ が $\mr_1$ の中の極小イデアルなら、(1) は
\[
        \mr_1=\mt+\mr\cap\mr_1 \qquad \hbox{§4, 3}
\]
を与え、$\mr_1$ が分解可能になってしまうからである。したがって $\mo$ は完全可約である。

第二に、\emph{極小条件を満たす根基なしの環は単位元をもつ。}

まず左単位元の存在を示すには、次を証明すれば十分である。\emph{右イデアル $\ma=c_1\mo\ne\mo$ が左単位元 $c_1$ をもつなら、より大きい長さをもつ右イデアル $c\mo$ が存在し、それも左単位元 $c$ をもつ。} すでに完全可約性、したがってすべての長さの有界性が示されているので、有限回で $\mo$ 自身の左単位元に達するからである。

したがって $\ma=c_1\mo$、$c_1^2=c_1$ とする。式
\[
        r=c_1r+(r-c_1r)
\]
により、上と同様に Peirce 分解 $\mo=c_1\mo+\mr$ が得られる。上と同様に $\mr$ の冪等元を $c_2$ とする。すなわち $c_2^2=c_2$ であり、$c_1$ は $\mr$ のすべての元を消すので $c_1c_2=0$ である。$e_1=c_1$、$e_2=c_2-c_2c_1$ とおく。このとき
\[
        e_1^2=e_1,
        \quad e_1e_2=0,
        \quad e_2e_1=0,
        \quad e_2c_2=e_2,
        \quad e_2\ne0,
        \quad e_2^2=e_2.
\]
上の注で述べたことにより、$c=e_1+e_2$ は環
\[
        c\mo=e_1\mo+e_2\mo
\]
の左単位元であり、$e_2^2=e_2\ne0$ であるから、この環は右イデアルとして $\ma$ より大きい長さをもつ。

構成された左単位元 $e$ が右単位元でもあることを示すため、左イデアルへの Peirce 分解
\[
        \mo=\mo e+\ml
\]
を行う。$\ml e=0$、$\ml=e\ml$ であるから、$\ml^2=\ml e\ml=0$ である。仮定により冪零左イデアルは存在しえないので、$\ml=0$。したがって $e$ は右単位元でもあり、結局、単位元である。

\emph{定理。右完全可約性と単位元の存在から両側完全可約性が従う：}
\[
        \mo=\md_1+\cdots+\md_s,
        \qquad \md_i \hbox{ は両側単純である。}
\]

前の証明と同じく、任意の極小両側イデアル $\ma$ が直和因子であることを示せば十分である。右イデアルとして、$\ma$ は直和因子である：
\[
        \mo=\ma+\mr=e_1\mo+e_2\mo.
\]
対応する左分解は
\[
        \mo=\mc+\ml=\mo e_1+\mo e_2
\]
である。すると
\[
        \mr\mc\subseteq \mr\ma\leq \mr\cap\ma=0,
        \qquad
        \ml\ma=\mo e_2e_1\mo=0,
\]
\[
        (\ma\ml)^2=\ma\ml\ma\ml=0,
        \qquad \hbox{したがって } \ma\ml=0;
\]
\[
        \mr=\mr\mo=(\mr\mc,\mr\ml)=\mr\ml,
        \qquad
        \ml=\mo\ml=(\ma\ml,\mr\ml)=\mr\ml.
\]
ゆえに $\mr=\ml$ であり、$\mr$ は両側である。したがって $\ma$ は両側直和因子である。ここから先の証明は、片側完全可約性の場合と同じである。

$\md_i$ は環としても両側単純である。$\md_i$ のすべてのイデアルは $\mo$ のイデアルだからである。これからそれらの構造を調べる。

\subsection*{§ 14. 単位元をもつ両側単純完全可約環}

$\mo$ をそのような環とする：
\[
        \mo=\mr_1+\cdots+\mr_n=e_1\mo+\cdots+e_n\mo,
        \qquad \mr_i\hbox{ は単純である。}
\]
このとき
\[
        \mo=\ml_1+\cdots+\ml_n=\mo e_1+\cdots+\mo e_n,
        \qquad \ml_i\hbox{ は直接分解不能である }\hbox{§9}.
\]
さらに $\ml_i\mr_i$ は非零両側イデアルである（$e_i^2$ が $\ml_i\mr_i$ に属するから）。したがってそれは $\mo$ である。

$\mA_{ik}=\mr_i\ml_k$ とおく。このとき
\[
        \mo=(\mr_1,\ldots,\mr_n)\cdot(\ml_1,\ldots,\ml_n)
        = (\ldots,\mA_{ij},\ldots)
\]
は直和である。実際、$0=\sum_{i,j}a_{ij}$ であれば、$a_{ij}$ は $\mr_i$ に属し、$\sum\mr_i$ は直和であるから、
\[
        0=\sum_j a_{ij},
\]
さらに $a_{ij}$ は $\ml_j$ に属するので、
\[
        0=a_{ij}
\]
である。

さらに $\mA_{ij}\mr_j=\mr_i\ml_j\mr_j=\mr_i\mo=\mr_i$ であるから、$\mA_{ij}\ne0$。$a_{ij}\ne0$ を $\mA_{ij}$ の元とする。このとき
\[
        a_{ij}\mr_j\subseteq\mr_i,
        \qquad
        a_{ij}\mr_j\ne0 \hbox{ である。なぜなら } a_{ij}e_j=a_{ij},
\]
ゆえに
\[
        a_{ij}\mr_j=\mr_i.
\]
$\mr_j$ の各 $r_j$ について
\[
        a_{ij}r_j=r_i
\]
とおくと、写像 $r_j\mapsto r_i$ は $\mr_j$ から $\mr_i$ への作用素準同型である。零へ写される元は $\mr_j$ に含まれるイデアルをなし、したがって零イデアルである。ゆえにこの写像は\emph{同型}である。したがって、

\emph{$\mo$ の一つの表示に現れる任意の二つの単純右イデアル $\mr_i$ と $\mr_j$ は作用素同型であり、その同型は $\mA_{ij}$ の元によって媒介される。} 異なる任意の二つの単純右イデアル $\mr,\mr'$ は、少なくとも一つの $\mo$ の表示にともに現れるので、そのような任意の二つのイデアルは同型である。

\emph{$\mr_j$ から $\mr_i$ へのすべての準同型は、$\mA_{ij}$ の元によって媒介される。}

\emph{証明。} $e_j\mapsto a_{ij}$ なら、$e_j^2\mapsto a_{ij}e_j$ であるから、$a_{ij}=a_{ij}e_j$ は $\mr_i\ml_j=\mA_{ij}$ に属する。さらに
\[
        \mr_j=e_j\mr_j\mapsto a_{ij}\mr_j=\mr_i.
\]

\emph{とくに、$\mr_i$ からそれ自身への準同型は $\mA_{ii}$ によって媒介される。}

$\mA_{ii}$ の二元の積は準同型の積に、和は和に対応する（定義は §1, 6 参照）。異なる $a_{ii}$ は異なる自己同型を与える。なぜなら $e_i$ は $a_{ii}e_i=a_{ii}$ に送られるからである。したがって環 $\mA_{ii}$ は $\mr_i$ の自己同型環と同型である。

しかし単純イデアルの自己同型環は斜体である（§2）。したがって $\mA_{ii}$ は斜体である。すべての $\mr_i$ は作用素同型であるから、それらの自己同型環は環同型である。すなわち $\mA_{ii}$ は $\mo$ によって環同型を除いて一意に定まる。現れうる異なる $\mA_{ii}$ は、単純右イデアルの抽象的に定義された自己同型環の、異なる具体的実現にすぎない。

\emph{主定理。単位元をもつ任意の両側単純完全可約環 $\mo$ は、斜体 $K$ 上の次数 $n$ の行列環と同型である。（斜体 $K$ は、$\mo$ の右イデアルの自己同型斜体と同型である。）}

\emph{証明。行列単位 $c_{ik}$ の構成。} $\Gamma_{11}$ を $\mr_1$ の恒等自己同型とし、$\Gamma_{i1}\mr_1=\mr_i$ である任意の同型を $\Gamma_{i1}$ とし、最後に
\[
        \Gamma_{ik}=\Gamma_{i1}\Gamma_{k1}^{-1}
\]
とおく。このとき一般に
\[
        \Gamma_{ik}\Gamma_{kl}=\Gamma_{il}
\]
である。$\Gamma_{ik}$ が $\mA_{ik}$ の元 $c_{ik}$ によって媒介されるなら、さらに
\[
        c_{ik}=e_i c_{ik}=c_{ik}e_k,
        \qquad
        c_{ik}c_{kl}e_l=c_{il}e_l,
\]
したがって
\[
        c_{ik}c_{kl}=c_{il},
        \qquad
        c_{ij}c_{kl}=c_{ij}e_j e_k c_{kl}=0\quad(j\ne k).
\]
ゆえに $c_{ik}$ は行列単位である（§7）。

\footnotetext{最近 B. L. van der Waerden は、抽象的な自己同型斜体を直接用いる、より簡単な証明を見いだした。この証明は彼の代数学の本（Grundlehren d. Math. Wiss., Berlin: Julius Springer）に掲載予定である。[1929年6月11日。]}

\emph{$K$ の構成。} 対応
\[
        a_{ij}=c_{i1}a_{11}c_{1i}
\]
は、$\mA_{11}$ の各 $a_{11}$ に、$\mA_{ii}$ の「共役元」$a_{ii}$ を対応させる。

この対応は環同型である。なぜなら、$a$ と $c$ に対応する同型 $\Delta,\Gamma$ は
\[
        \Delta_{ii}=\Gamma_{i1}\Delta_{11}\Gamma_{i1}^{-1}
\]
によって結ばれており、この関係は環同型を表すものとしてよく知られているからである。いま元
\[
        \alpha=a_{11}+\cdots+a_{nn}=a_{11}+c_{21}a_{11}c_{12}+\cdots+c_{n1}a_{11}c_{1n}
\]
を作る。各 $a_{11}$ はちょうど一つの $\alpha$ を定める（和が直和であるため一意である）。$\alpha$ の全体は斜体 $K\sim\mA_{11}$ である。実際、
\[
        \alpha=a_{11}+\cdots+a_{nn},
        \qquad
        \beta=b_{11}+\cdots+b_{nn}
\]
から
\[
        \alpha+\beta=(a_{11}+b_{11})+\cdots+(a_{nn}+b_{nn}),
        \qquad
        \alpha\beta=a_{11}b_{11}+\cdots+a_{nn}b_{nn}
\]
を得るので、対応 $a_{11}\mapsto\alpha$ は同型である。

さらに
\[
        c_{ik}\alpha=c_{ik}a_{kk}=c_{ik}c_{k1}a_{11}c_{1k}=c_{i1}a_{11}c_{1k},
\]
\[
        \alpha c_{ik}=a_{ii}c_{ik}=c_{i1}a_{11}c_{1i}c_{ik}=c_{i1}a_{11}c_{1k},
\]
であるから、任意の $\alpha$ はすべての $c_{ik}$ と可換である。

最後に、同じ式から
\[
        c_{ik}K=c_{ik}\mA_{kk}=c_{ik}\mr_k\ml_k=\mr_i\ml_k=\mA_{ik},
\]
したがって
\[
        \mo=\sum c_{ik}K
\]
が得られる。

\emph{対応。} $\mo$ の各元を $\sum c_{ik}\alpha_{ik}$ の形に書き、その元に行列 $(\alpha_{ik})$ を対応させると、この対応は行列乗法の定義から直ちに同型であることが従う。これで主定理が証明された。

\emph{逆。斜体 $A$ 上の次数 $n$ の行列環は両側単純かつ完全可約である。$A$ は右イデアルの自己同型環と同型であり、$Ae$ は前の定理で構成された $K$ である。}

\emph{証明。} 行列単位（§7）を $c_{ik}$ とする。
\[
        \mr_i=\sum_k c_{ik}A
\]
とおく。このとき明らかに $\mo=\mr_1+\cdots+\mr_n$ である。$\mr_i$ は単純右イデアルである。任意の非零元 $\sum c_{ik}\alpha_k$ が $\mr_i$ 全体を生成するからである。実際、$\alpha_j\ne0$ なら
\[
        \Bigl(\sum c_{ik}\alpha_k\Bigr)\cdot(c_{jl}\alpha_j^{-1})=c_{il},
\]
であり、$\sum c_{il}\beta_l$ は $\mr_i$ のすべての元を走る。

$\mr_i$ は作用素同型である：$\mr_i=c_{ik}\mr_k$。

したがって $\mo$ は両側分解不能である。ゆえに、§13 の定理とすでに証明した完全可約性により、また単位元が存在するので、$\mo$ は両側単純である。このことは、任意の $\sum\sum c_{ik}\alpha_{ik}\ne0$ が両側イデアル $\mo$ 全体を生成することからもわかる。

さらに
\[
        \mr_i=c_{ii}\mo,
        \qquad
        \ml_i=\mo c_{ii},
\]
\[
        e=\sum c_{ii},
        \qquad c_{ii}=e_i,
\]
\[
        \mA_{ii}=\mr_i\cap\ml_i=c_{ii}A\sim A.
\]
最後に、$a_{11}=c_{11}\lambda$ とおくと、
\[
        \alpha=a_{11}+c_{21}a_{11}c_{12}+\cdots
        =c_{11}\lambda+c_{22}\lambda+\cdots=e\lambda,
\]
が得られ、主張はすべて証明される。

\emph{$\mo$ の中心は $K$ の中心であり、したがって体である。}

\emph{証明。} $\mZ(K)$ はすべての $x$ と可換なすべての $\zeta$ からなる。しかしこれらはすべての $c_{ik}$ とも可換であり、したがってすべての和 $\sum c_{ik}a_{ik}$ と可換である。ゆえに
\[
        \mZ(K)\subseteq\mZ(\mo).
\]
逆に、$z$ を $\mZ(\mo)$ の元とする：
\[
        z=\sum_\mu\sum_\nu c_{\mu\nu}\gamma_{\mu\nu}.
\]
このとき
\[
        zc_{ik}=c_{ik}z,
        \qquad
        \sum_\mu c_{\mu k}\gamma_{\mu i}=\sum_\nu c_{i\nu}\gamma_{k\nu},
\]
したがって
\[
        \gamma_{\mu i}=0\ (\mu\ne i),
        \qquad
        c_{ik}\gamma_{ii}=c_{ik}\gamma_{kk},
        \qquad
        \gamma_{ii}=\gamma_{kk}=\gamma;
\]
\[
        z=\sum c_{ii}\gamma=e\gamma=\gamma.
\]
したがって $z$ は $K$ に属し、すべての $x$ と可換であるから、$z$ は $\mZ(K)$ に属する。

行列環はもちろん左完全可約でもあり、単位元をもつ任意の右完全可約環は行列環の直和であるから、次が従う：

\emph{単位元をもつ任意の右完全可約環は左完全可約でもあり、逆も成り立つ。}

また、\emph{単位元をもつ完全可約環の中心は、両側単純行列環に対応する可換体の直和である。}

次の定理が残る：

\emph{任意の二つの異なる分解 $\mo=\sum c_{ik}K$、$\mo=\sum c'_{ik}K'$ は、内部自己同型 $\alpha'\mapsto x^{-1}\alpha'x$、$c'_{ik}\mapsto x^{-1}c_{ik}x$ によって互いに移される。}

\emph{証明。} 
\[
        \mo=\mr_1+\cdots+\mr_n=\mr_1'+\cdots+\mr_n'
\]
とし、$a,b$ を、イデアル $\mr_1,\mr_1'$ の二つの逆向きの同型を媒介する元とする：
\[
        \mr_1=a\mr_1',
        \qquad
        \mr_1'=b\mr_1,
        \qquad
        ab=c_{11},
        \qquad
        ba=c'_{11}.
\]
次をおく：
\[
        x=\sum_i c_{i1}ac'_{1i},
        \qquad
        y=\sum_i c'_{i1}bc_{1i}.
\]
すると
\[
        xy=e,
        \qquad
        y=x^{-1},
        \qquad
        x^{-1}c_{ik}x=c'_{ik},
        \qquad
        x^{-1}\alpha x=\alpha'.
\]
\end{document}
